\begin{theorem}[Cauchy subsequence dependency route]
\label{thm:cauchy-acceleration-cauchy-subsequence-route}
Every $\CauchyAccelerationUp$ consumer that exports an accelerated $\RegSeqRatUp$ readback or terminal $\RealUp$ seal must first consume the rate-dominance surface of \autoref{ch:concrete-instances-cauchy-rate-dominance-namecert} and the fast-modulus surface of \autoref{ch:concrete-instances-fast-cauchy-modulus-namecert}. The route is finite: $\CauchyRateDominanceUp$ bounds the source schedule, $\FastCauchyModulusUp$ supplies the inspected acceleration window, and the displayed $\StreamNameUp$, $\DyadicRatCoreUp$, $\RegSeqRatUp$, $\RealUp$, transport, replay, provenance, and local $\NameCert$ rows carry the resulting readback. No consumer may replace this route by a chosen subsequence, quotient representative, ambient completion row, or host equality of Real names.
\end{theorem}
\begin{proof}
Begin with the accepted carrier of \autoref{def:cauchy-acceleration-carrier}. Its source stream row $S$, tolerance schedule $T$, cofinal ledger $\alpha$, regular rational row $R$, and Real seal $E$ are the only analytic rows available before transport, replay, provenance, and local naming.

The rate-dominance chapter supplies the comparison row that bounds which source schedule entries may be read by $\alpha$. The fast-modulus chapter supplies the inspected window that turns those bounded reads into the accelerated tolerance route. Reading either row after the $\RegSeqRatUp$ readback would hide the schedule obligation rather than consume it.

Apply the rectangular window route of \autoref{thm:cauchy-acceleration-rectangular-window-route}. It keeps the $\StreamNameUp$ and $\DyadicRatCoreUp$ rows in front of $R$ and $E$, so the accelerated readback remains a displayed finite route. The carrier has no coordinate for a selected subsequence, quotient representative, ambient completion proof, or host Real equality, and the dependency route cannot export one.
\end{proof}
